\begin{table}[t]
\centering
\small
\begin{tabular}{llll}
\toprule
Group & Features & Why it should help & Hypothesis \\
\midrule
Transforms & log1p of skewed numerics (skew > 0.75, decided per fit); sin/cos of MoSold & makes area–price relations linear in logs; month becomes cyclic so Dec and Jan are neighbours & H1 \\
Interactions & Qual\_x\_Area, TotalQualCond & lets the quality effect grow with size (luxury premium); one combined condition score & H2, H5 \\
Aggregates & TotalArea, TotalBaths, TotalPorchArea & buyers value total usable space; one strong feature replaces several correlated parts & H1, H7 \\
Domain-derived & HouseAge, RemodAge, GarageAge; AvgRoomSize, LotUtilization & age at sale is what a buyer sees, not the calendar year; room size and lot use measure spaciousness & H1, H3 \\
Presence flags & HasBasement, HasFireplace, HasPool, HasGarage, IsRemod & capture the price step of having an amenity, separate from its size & H6 \\
\bottomrule
\end{tabular}
\caption{Engineered feature groups, the rationale for each and the hypothesis it addresses (§4).}
\label{tab:feature_groups}
\end{table}
